%! Comparing Studies and Choosing a Method — Unit 1, Lesson 1.7 — Homework
\documentclass[10pt]{article}
\usepackage{saar-article}
\usepackage{saar-boxes}
\newcommand{\TallMath}[1]{$\displaystyle #1\rule[-1.4em]{0pt}{3.2em}$}

\begin{document}

\pageheader{Unit 1, Lesson 1.7}{Homework}
% NAMESTRIP: no \namedateperiod here. The student writes their name once, on the
% cover this component is stapled behind. See references/conventions.md.

\vspace{0.15in}

\boxguard[24]
\begin{scenariobox}[Millbrook Public Library --- The Third Study]{cerulean}
\small
Millbrook has $1{,}200$ card holders and has now studied its Summer Reading
Program twice. \textbf{Last summer it watched:} of the card holders who had
signed themselves up, $80\%$ read at least $10$ books against $32\%$ of those
who had not --- a $48$-point gap. \textbf{This summer it assigned:} $200$
volunteers split $100$/$100$ by a generator, $65\%$ against $35\%$ --- a
$30$-point gap.

\vspace{2pt}
\textbf{This fall the director runs a third study.} Does a weekly evening book
club raise how many card holders read at least $10$ books in a semester? She
takes \textbf{$300$ volunteers} and lets a generator send \textbf{$150$ to the
book club} and \textbf{$150$ to no club}, keeping borrowing limits the same for
everyone. By December, \textbf{$96$} of the club group and \textbf{$60$} of the
no-club group had read at least $10$ books.
\end{scenariobox}

\vspace{0.08in}

\boxguard
\begin{notesbox}{Practice}
\small
\begin{enumerate}[leftmargin=1.6em, itemsep=8pt, topsep=3pt]

  \item Sort the two earlier studies.

        \textbf{(a)} Last summer's is a(n) \blank{3.2cm} \hfill This summer's is
        a(n) \blank{3.2cm}

        \textbf{(b)} Who decided which group each card holder was in last
        summer? \blank{5.4cm}

  \item Compute the book-club rates.

        \textbf{(a)} What percent of the $150$ in the club read at least $10$
        books?

        \begin{work}
          \dfrac{96}{150} &= 0.64 \\
                          &= 64\%
        \end{work}

        \textbf{(b)} What percent of the $150$ with no club did?

        \begin{work}
          \dfrac{60}{150} &= 0.40 \\
                          &= 40\%
        \end{work}

        \textbf{(c)} The club group is higher by \blank{2cm} percentage points.

  \item Now use all $300$ volunteers.

        \textbf{(a)} How many read at least $10$ books, and what percent is
        that?

        \begin{work}
          96 + 60 &= 156 \text{ card holders} \\
          \dfrac{156}{300} &= 0.52 = 52\%
        \end{work}

        \textbf{(b)} About how many of all $1{,}200$ card holders does that
        predict?

        \begin{work}
          0.52 \times 1200 &= 624 \text{ card holders}
        \end{work}

  \item Fill in the plan the director followed for the book-club study.

        \begin{center}
        \renewcommand{\arraystretch}{1.4}
        \begin{tabularx}{0.93\linewidth}{@{} p{2.4cm} X @{}}
        \toprule
        \textbf{Stage} & \textbf{The book-club study} \\
        \midrule
        Formulate & Does a weekly evening book club raise how many card holders
          read at least $10$ books? \\
        Collect & \blank{7.2cm} \\
        Represent & \blank{7.2cm} \\
        Analyze \& report & \blank{7.2cm} \\
        \bottomrule
        \end{tabularx}
        \end{center}

  \tcbbreak   % keep item 5's prompt with its table — mirrored in homework_key
  \item The board sends three more questions. For each, decide whether the
        director could assign the groups.

        \begin{center}
        \renewcommand{\arraystretch}{1.45}
        \begin{tabularx}{0.97\linewidth}{@{} X c p{4cm} @{}}
        \toprule
        \textbf{The question} & \textbf{Assign?} & \textbf{If not, why --- and
          what instead?} \\
        \midrule
        Does a $\$5$ late fee reduce overdue books? & \blank{1.1cm}
          & \blank{3.6cm} \\
        Do card holders who never finished high school borrow fewer books?
          & \blank{1.1cm} & \blank{3.6cm} \\
        Do children who were read to as toddlers read more at age $12$?
          & \blank{1.1cm} & \blank{3.6cm} \\
        \bottomrule
        \end{tabularx}
        \end{center}

  \item Name the collection method that fits each job best.

        \begin{center}
        \renewcommand{\arraystretch}{1.45}
        \begin{tabularx}{0.97\linewidth}{@{} X p{3.6cm} @{}}
        \toprule
        \textbf{What the library needs} & \textbf{Best method} \\
        \midrule
        Which titles were borrowed most over the past five years
          & \blank{3.2cm} \\
        How many of the $1{,}200$ card holders want Sunday hours
          & \blank{3.2cm} \\
        Why one long-time member stopped coming in, in her own words
          & \blank{3.2cm} \\
        How many visitors actually use the new self-checkout machine
          & \blank{3.2cm} \\
        Six teens talking together about what the teen room needs
          & \blank{3.2cm} \\
        \bottomrule
        \end{tabularx}
        \end{center}

  \item Millbrook now has three results: a $48$-point gap, a $30$-point gap, and
        a $24$-point gap. Which one is the \emph{weakest} evidence that the
        library's programs work, and why is it the one with the biggest number?
        \writelines{2}

  \item A local paper runs the headline \emph{``Millbrook study proves the
        Summer Reading Program makes children read more,''} citing last
        summer's study. Rewrite the headline so that it is honest, and say what
        the library would have to do to earn the original wording.
        \writelines{2}

\end{enumerate}
\end{notesbox}

\vspace{0.08in}

\boxguard
\begin{extensionbox}
\small
Go back to the study you proposed for Lesson 1.6's extension and finish planning
it:

\begin{enumerate}[leftmargin=1.6em, itemsep=3pt, topsep=3pt]
  \item Can it actually be run as an experiment, or does something stop you from
        assigning? Name which of the three reasons applies if so.
  \item Write your plan as the four stages of the statistical cycle, one line
        each.
  \item Which of the five collection methods will you use, and why that one?
  \item Write, in advance, the strongest sentence your design will let you say
        --- before you know how the numbers come out.
\end{enumerate}
\end{extensionbox}

\vspace{0.08in}

\begin{remindbox}
\small
\textbf{Next class (Lesson 1.8):} you run the study. Working from a question your
class chooses, you will plan the whole cycle, pick a sampling technique and a
collection method, collect real data from real classmates, and report a
conclusion your design actually earns.
\end{remindbox}

\end{document}
